\section{Conclusions}
\label{sec:conclusions}

The long-horizon error of autoregressive surrogates was formulated as a forced finite-time problem and decomposed exactly into the magnitude of the injected defects, the worst-case gain of the tail propagators, the efficiency with which the defects use that gain, the cancellation in the vector sum and a nonlinear remainder. Null models for the efficiency and for the vector composition, together with an exact variance attribution of the log-error, separate the properties of the learned dynamics from those of the learned errors.

For residual MLPs trained on Lorenz--96 with three optimization budgets and for FNOs trained on Kuramoto--Sivashinsky, the worst-case first-order estimate exceeds the realized error by a factor of five to seven at about one Lyapunov time. The defects realize the fraction of the worst-case gain expected for random vectors with their own covariance, but add about four times more coherently than independent ones, because consecutive defects are correlated along the trajectory. The defect magnitude explains which forecasts degrade first, whereas the worst-case gain, the coherence and, for the neural operators, the orientation of the defects explain which forecasts are worse beyond about half a Lyapunov time. These conclusions hold over a sixtyfold range of one-step accuracy, and the first-order description is accurate while the error remains below about 15 to 20\% of the saturation level.

Since the realized error depends on the actual defect field, the disagreement between independently trained surrogates provides a reference-free estimate of it, which ranks individual forecasts up to about one Lyapunov time and, for the converged Lorenz--96 surrogates, predicts the error in magnitude. Multistep fine-tuning reduces the error of under-trained surrogates by damping their tangent dynamics, but it has no consistent effect on surrogates that reproduce the reference tangent dynamics, so that improvements from multistep objectives should be judged against converged one-step baselines. Natural extensions are the use of the factorization with matrix-free estimates for high-dimensional neural operators and weather models, and training objectives that act on the coherence of the transported defects.
